% SOURCE_STATEMENT: "1. There exists an explicit construction of Atiyah [1] and Page [10] for the Green's function on certain ALE\n\nand ALF spaces. Page's construction comes from physics and constructs by hand the Green's functions for both AkALE and AkALF. Atiyah has a formalism which in theory allows one to construct the Green's function on any self-dual gravitational instanton - this formalism involves Serre classes on the corresponding twistor spaces. In the paper he constructs only the Green's function in the AkALE case. (According to Etesi the use of Atiyah's Serre-class formalism is significantly more difficult for the ALF case and has not yet been carried out. The main problem he says is that one cannot find a nice compacification for the twistor space of an ALF self-dual space. Carrying out this construction would allow us to construct L2\n\nharmonic forms over the new DkALF gravitational instantons constructed by Cherkis and Hitchin.)\n\nQuestion: Is there a way to generalize this formalism to 8 dimensional hyperk¨ ahler manifolds?\n\nReason: In the non-compact 4-dimensional case or K3 case, the coefficient of the leading order term in the Green's function gave information about the values of the parameters (e.g. sizes of the cycles, metric). 2. The same story for quaternionic manifolds. How to modify Atiyah's formalism?\n\nReason: There is a reasonable twistor theory for these manifolds. 3. Is it true that any 4-dimensional hyperk¨ ahler manifold (non-compact, with reasonable decay at infinity) can be deformed to a quaternionic one?\n\nFor example: HK has Ric = 0, while quaternionic manifolds satisfy Ric = λg with λ a constant (non-zero cosmological constant). One would like to obtain any HK as a limit of quaternionic spaces with λ → 0.\n\nExample: R4 is the limit of HP 1 (positive λ's) and HH 1 (negative λ's).\n\n2 March 18, 2004 - S. Cherkis as moderator.\n\nQuestion (R. Mazzeo):\n\nWhat are the asymptotics at infinity of the multi-monopole moduli spaces?\n\nAnswer: (Cherkis) They should behave like a N-body problem.\n\nQuestions (K. Lee)\n\nFor these questions, R4\n\n> B\n\ndenotes the non-commutative R4. Here B = Bμν dx μdx ν is a 2-form which gives the non-commutative structure. Then, [ xμ, x ν ] = Bμν. We define non-commutative instantons using the Nahm transform of instanton data\n\nμ1 = [ T1, T 2] − [T3, T 4] − j∗j = B12,\n\nplus two other similar equations (see the work Nekrasov-Schwarz on non-commutative instantons [9]). Let\n\nMk,n B be the moduli space of centered k-instantons on R4\n\n> B\n\nwith gauge group G = U (n). (The fact that the instantons are centered means we consider only instantons with center of mass at the origin) 1. What is the instanton number in the non-commutative set-up?\n\nMotivation: The space Mk,n B is singular for B = 0 but non-singular for B 6 = 0. One would like to understand the limit as B → 0. 2. Does there exists a unique harmonic form in the middle dimension of Mk,n B?\n\nNote: The answer to this question is Yes for the following cases: (i) 1-instantons and any G;(ii) 2-instantons and G = U (1) (see the work of Lee et al). 33. Conjecture: The moduli space of centered 1-instantons on R3 × S1\n\n> B\n\nwith \"given holonomy around S1\n\n> ∞\n\n\"(holonomy is not in the center of G = U (n)) has exactly n L 2-harmonic forms in middle dimension."
% Preparation and final self-audit: CONTRIBUTING.md prompts C and D; source/audit artifacts retained in the local physics_statements_2026_10_11 batch.
% !TeX program = lualatex
\documentclass[11pt]{article}
\usepackage[margin=25mm]{geometry}
\usepackage{fontspec}
\setmainfont{Latin Modern Roman}
\usepackage{amsmath,amssymb,mathtools}
\usepackage{unicode-math}
\setmathfont{Latin Modern Math}
\usepackage{xurl,hyperref}
\hypersetup{colorlinks=true,urlcolor=blue}
\setcounter{secnumdepth}{0}
\setlength{\parindent}{0pt}
\setlength{\parskip}{5pt}
\setlength{\emergencystretch}{3em}
\newcommand{\sref}[2]{\hyperlink{source:#1:#2}{[S#2]}}
\newcommand{\cataloguescope}[1]{\paragraph{Recorded target.}#1}
\newcommand{\R}{\mathbb R}
\newcommand{\C}{\mathbb C}
\newcommand{\N}{\mathbb N}
\newcommand{\Z}{\mathbb Z}
\newcommand{\Q}{\mathbb Q}
\newcommand{\F}{\mathbb F}
\newcommand{\eps}{\varepsilon}
\begin{document}
\section{A synchronized heat-kernel construction on eight-dimensional hyperkähler products}\label{statement}

\subsection{Definitions and mathematical statement}

\paragraph{Shared notation.}$\mathbb N=\{1,2,\ldots\}$; the standard number fields have their usual meanings. All additional conventions are stated locally.


\paragraph{Definitions and notation.}A hyper-Kähler metric has three parallel complex structures satisfying the quaternionic relations. In four dimensions a gravitational instanton is a complete noncompact Ricci-flat geometry in the source’s self-dual setting. ALE and ALF describe different asymptotic geometries; ALF has an asymptotic circle direction over a three-dimensional base. For a Riemannian manifold $(M,g)$, $L^2$ norms use the metric volume and pointwise metric on forms. Reduced $L^2$ cohomology is $\ker d/\overline{\operatorname{im}d}$ for the specified closed $L^2$ complex. On a complete manifold its harmonic representatives satisfy $d\alpha=d^*\alpha=0$ and $\int_M|\alpha|^2<\infty$. A Green function is a distributional inverse kernel for a specified elliptic operator and boundary behavior. The twistor space of a hyper-Kähler manifold encodes its sphere of complex structures. Quaternionic-Kähler geometry in this context has Einstein metric with nonzero scalar-curvature parameter. Noncommutative instanton moduli in the source use commutators $[x_\mu,x_\nu]=B_{\mu\nu}$ and moment-map equations deformed by $B$. Centered means the center-of-mass parameter is removed.

\paragraph{Catalogue wording retained for traceability.}
Generalize Atiyah’s Serre-class Green-function formalism from self-dual four-manifolds to eight-dimensional hyper-Kähler manifolds, and formulate its analogue for quaternionic manifolds. Determine whether every four-dimensional hyper-Kähler manifold with suitable noncompact decay is a limit as $\lambda\to0$ of quaternionic Einstein metrics $\operatorname{Ric}=\lambda g$. Determine asymptotic geometry of multi-monopole moduli spaces, motivated by N-body behavior. For centered $k$-instanton moduli $M_B^{k,n}$ on noncommutative $\mathbb R_B^4$ with gauge group $U(n)$ and deformed moment-map equations, define instanton number and determine whether there is a unique harmonic form in middle dimension. The source records affirmative cases for one instanton and for two instantons with $U(1)$. Prove or disprove that centered one-instanton moduli on $\mathbb R^3\times S_B^1$ with specified noncentral holonomy at infinity have exactly $n$ $L^2$-harmonic middle-dimensional forms.
\paragraph{Scope and source qualification.}All questions in this canonical excerpt are retained; the product heat-kernel construction suggested by its title is only a possible subcase. The source’s informal decay class, harmonic-form integrability and noncommutative conventions require specification. Historical special-case statements are not independently verified. The extracted background remarks spill into separate subsequent massless-monopole, Singer and wave-trace questions; these are not silently reassigned to this canonical Cherkis/Mazzeo/Lee excerpt. The catalogue title/status was paired with a machine-generated research attempt; that attempt is not peer-reviewed evidence of settlement. The statement covers the original source questions rather than only the title’s subcase.
\subsection{Short English statement}
Generalize Green-function twistor constructions, study hyper-Kähler limits of quaternionic metrics, and count harmonic forms on noncommutative instanton moduli.
\subsection{Sources}
\begin{itemize}
\item[\hypertarget{source:statement:1}{S1}] ULAM.AI and the UnsolvedMath Contributors, UnsolvedMath, catalogue entry. Catalogue curation under CC BY 4.0; definitions and English exposition added locally.
\url{https://huggingface.co/datasets/ulamai/UnsolvedMath}.
\item[\hypertarget{source:statement:2}{S2}] American Institute of Mathematics, $L^2$ harmonic forms in geometry and string theory; notes by Anda Degeratu and Mark Haskins, March 2004, 17–18 March 2004, Cherkis, Mazzeo and Lee questions.
\url{https://aimath.org/WWN/l2harmonic/l2harmonic.pdf}.
\end{itemize}
\label{end:statement}
\end{document}
